
This study introduced C\_sem, a training-free SBERT-based metric for measuring semantic accommodation in human-AI conversation, and applied it to 155,362 turn pairs from the Anthropic HH-RLHF helpfulness dataset. Three findings are supported by the data: (1) human follow-up turns exhibit substantial interaction-specific semantic similarity to their AI partner turns (C\_sem = 0.329, d = 1.998 vs. random, d = 0.417 vs. KNN topic-matched); (2) this similarity increases monotonically with RLHF tier quality (J-T p = 0.001, d = 0.183--0.254 for T1 to T3, replicated across three SBERT architectures); and (3) human-to-AI accommodation exceeds AI-to-human accommodation in all nine tier-by-model conditions (d = 0.061--0.405).

Two mechanism hypotheses are falsified. Within-conversation quality discrimination yields a reversed signal: human follow-ups are more similar to rejected than chosen AI responses (delta < 0 in 25 of 27 conditions, d up to -0.738). A politeness-style mediation proxy has no predictive power (beta approximately 0, p approximately 0.99). The tier-accommodation association is therefore characterized as a population-level pattern rather than a within-conversation perceptual effect. The cross-sectional design precludes causal claims; longitudinal within-user studies would be needed to establish whether RLHF quality causally influences human semantic behavior.
